% Source: https://artofproblemsolving.com/community/c3h1336697_amc_8_2005_15
% Stable identity: aops-topic-1336697-q15
Let the two equal sides be $a$ and the third side $b$, so $2a+b=23$. Then $b=23-2a$ is odd (no triangle is equilateral, since $3\nmid23$). The triangle inequality requires $b<2a$, i.e. $23-2a<2a$, so $a\ge6$; and $b\ge1$ gives $a\le11$. So $a=6,7,8,9,10,11$ with $b=11,9,7,5,3,1$: there are $\boxed{6}$ triangles, choice $\textbf{(C)}$.

Counting every odd base (which would give $11$, choice (E)) ignores the triangle inequality: for example $a=5$, $b=13$ is impossible because $5+5<13$. Also, the equal sides need not be even; only $b$ must be odd.

The two extreme cases, $(6,6,11)$ and $(11,11,1)$, drawn to scale:
\begin{asy}
usepackage("amsmath");
size(6cm);
defaultpen(fontsize(10pt));
real s=1/11; pair B2=(0,0),C2=(11*s,0),A2=(5.5*s,sqrt(36-30.25)*s);
draw(B2--C2--A2--cycle,linewidth(0.9));
label("$6$",(B2+A2)/2,NW); label("$6$",(C2+A2)/2,NE); label("$11$",(B2+C2)/2,S);
pair o=(1.3,0); real t=1/11;
draw(o--o+(t,0)--o+(0.5t,sqrt(121-0.25)*t)--cycle,linewidth(0.9));
label("$11$",o+(0.25t,5.5t),W); label("$11$",o+(0.75t,5.5t),E); label("$1$",o+(0.5t,0),S);
label("\mbox{}",(min(currentpicture,user=true)+max(currentpicture,user=true))/2,black); // keeps the text after the figure black when LaTeX typesets the labels inline
\end{asy}
